\documentclass[10pt,a4paper]{amsart}
\usepackage[margin=19mm]{geometry}
\usepackage{amsmath,amssymb,mathtools}
\usepackage[scr=rsfs,cal=euler]{mathalfa}
\usepackage{xfrac}
\usepackage[unicode,hidelinks,bookmarks=false]{hyperref}
\newcommand{\ie}{i.e.\ }
\newcommand{\cf}{cf.\ }
\newcommand{\resp}{respectively\ }
\newcommand{\Subsection}{\subsection}
\providecommand{\nobreakdash}{\nobreak\hbox{-}}
\setlength{\emergencystretch}{2em}
\multlinegap0pt
\usepackage{bm}
\usepackage{bbm}
\usepackage{dsfont}
\usepackage{xspace}
\usepackage{cancel}
\usepackage{mathtools}
\usepackage{amsthm}
\makeatletter
\@ifundefined{theorem}{\newtheorem{theorem}{Theorem}}{}
\@ifundefined{proposition}{\newtheorem{proposition}[theorem]{Proposition}}{}
\makeatother
\newcommand{\E}{\mathbb{E}}
\newcommand{\Lgen}{\mathscr{L}}
\newcommand{\dd}{\,\mathrm{d}}
\title[Delay-Penalty Comparison for Sequential Testing and Quickest]{Delay-Penalty Comparison for Sequential Testing and Quickest Detection in State-Dependent Diffusion Models}
\author{Ye Liang}
\date{Source: arXiv:2606.23962v1 (CC BY 4.0)}
\begin{document}
\maketitle

\noindent\emph{Source and licence.} Delay-Penalty Comparison for Sequential Testing and Quickest Detection in State-Dependent Diffusion Models. Version arXiv:2606.23962v1 (CC BY 4.0), distributed under the Creative Commons Attribution 4.0 International licence, as recorded in the arXiv licence metadata for this exact version and shown on the abstract page https://arxiv.org/abs/2606.23962v1. Full rights notice: licenses/arxiv-2606.23962-CC-BY-4.0.txt.\par\smallskip

\noindent\textit{Editorial scope.} Counted unit: the Proposition whose source label is \texttt{prop:verif} in the article (source0.tex lines 723--740, source label \texttt{prop:verif}), delivered together with the article's own complete original proof of it (source0.tex lines 742--764, one \texttt{proof} environment of 23 lines). No mathematics is written, repaired or re-derived: statement and proof are the article's own text line for line. Required context retained uncounted: none. Every definition, display and label of those lines is retained unchanged. Omitted from this package and not redistributed: 1--722, 741--741, 765--905. Nothing inside a retained excerpt is omitted or altered: every retained body is delivered exactly as the article's lines are written, so no omission receipt of any kind accompanies this package. Citations: the retained excerpts carry no citation command at all, so no bibliography entry of the article is transcribed and no reference is redistributed. Reference adaptation: none was needed and none was made; every reference inside the delivered unit resolves inside it, so no unresolved reference or citation is produced. Printed slips kept literal: no printed word, symbol, hypothesis or number of the counted statement or of its proof was altered. Layout and class: the article's own class is \texttt{sn-jnl} with packages including amsmath, amssymb, amsthm, mathtools, graphicx, multirow, amsfonts, mathrsfs, appendix, xcolor, textcomp, manyfoot, booktabs, algorithm; the wrapper is the lane's pinned \texttt{amsart} editorial wrapper at 10pt on A4, which restates the theorem-like environments the retained text needs and adds the article's own macro definitions that the retained text needs (\texttt{\textbackslash E}, \texttt{\textbackslash Lgen}, \texttt{\textbackslash dd}), with extra wrapper packages (bm, bbm, dsfont, xspace, cancel, mathtools). The delivered numbering is the unit's own. Assets: no retained span contains a figure environment, a graphics-inclusion command, a PDF-page inclusion, a table or a TikZ picture; no image file is copied into this package, the package declares no asset dependency and no image file is redistributed with it. Licence: version arXiv:2606.23962v1 of this article is distributed under the Creative Commons Attribution 4.0 International licence, as recorded in the arXiv licence metadata for this exact version and shown on the abstract page https://arxiv.org/abs/2606.23962v1. A full rights notice is delivered at licenses/arxiv-2606.23962-CC-BY-4.0.txt. Characters: every retained excerpt is pure ASCII, so the delivered engine, XeLaTeX with its default Latin Modern text font, reports no missing character.
\begin{proposition}[Verification]\label{prop:verif}
Let $\widehat V$ be continuous on $\mathcal S$, of polynomial growth, $C^1$ across
$\partial\widehat{\mathcal C}$, and $C^2$ on the interiors of
$\widehat{\mathcal C}=\{\widehat V<G\}$ and $\widehat{\mathcal D}=\{\widehat V=G\}$,
and suppose
\[
  \Lgen\widehat V+f\ge0\ \text{on }\mathcal S,\qquad
  \widehat V\le G\ \text{on }\mathcal S,\qquad
  \Lgen\widehat V+f=0\ \text{on }\widehat{\mathcal C}.
\]
Suppose moreover that for every admissible $\tau$ the local martingale
$M_{\,\cdot\,}$ in \eqref{eq:ito-verif} below, stopped along a localizing sequence
$\tau_n\uparrow\infty$, is uniformly integrable in the limit (e.g.\ a
square-integrability or sublinear-growth condition on
$\nabla\widehat V\cdot\sigma(Y)$). Then $\widehat V=V$, and
$\tau^\ast=\inf\{t:Y_t\in\widehat{\mathcal D}\}$ is optimal whenever
$\E_y[\tau^\ast]<\infty$.
\end{proposition}

\begin{proof}
For an admissible $\tau$ and a localizing sequence $\tau_n\uparrow\infty$, the
It\^o--Tanaka formula applied to $\widehat V(Y_{\cdot})$ gives
\begin{equation}\label{eq:ito-verif}
  \widehat V(Y_{\tau\wedge\tau_n})+\int_0^{\tau\wedge\tau_n} f(Y_s)\dd s
   =\widehat V(y)+\int_0^{\tau\wedge\tau_n}(\Lgen\widehat V+f)(Y_s)\dd s
   +M_{\tau\wedge\tau_n},
\end{equation}
where $M$ is a local martingale. The $C^1$ (smooth-fit) hypothesis across
$\partial\widehat{\mathcal C}$ ensures that the local-time term on the free
boundary, which would otherwise appear because $\Lgen\widehat V$ has a jump in its
second derivatives there, vanishes; where only continuous fit holds, the
local-time term is nonnegative and is retained in the inequality below without
affecting its direction. Since $\Lgen\widehat V+f\ge0$, taking expectations and
letting $n\to\infty$ (using the growth and uniform-integrability hypotheses so
that $\E_y[M_{\tau\wedge\tau_n}]\to0$) gives
$\widehat V(y)\le\E_y[\int_0^\tau f\dd s+\widehat V(Y_\tau)]
\le\E_y[\int_0^\tau f\dd s+G(Y_\tau)]$, hence $\widehat V\le V$. For
$\tau=\tau^\ast$ the integrand $\Lgen\widehat V+f$ vanishes on
$\widehat{\mathcal C}$ and $\widehat V(Y_{\tau^\ast})=G(Y_{\tau^\ast})$, so the
inequalities are equalities and $\widehat V(y)=\E_y[\int_0^{\tau^\ast}f\dd
s+G(Y_{\tau^\ast})]\ge V(y)$. Thus $\widehat V=V$ and $\tau^\ast$ is optimal.
\end{proof}

\end{document}
